\documentclass[12pt]{article}

\usepackage[margin=1.25in]{geometry}
\usepackage{booktabs}
\usepackage{array}
\usepackage{longtable}
\usepackage{hyperref}
\usepackage{setspace}
\usepackage{parskip}

\onehalfspacing

\title{The Independent Treasury: Origins, Structure, and Operations}
\author{}
\date{March 2026}

\begin{document}

\maketitle

\tableofcontents
\bigskip

%-----------------------------------------------------------------------
\section{Background: The Bank Wars (1832--1836)}
%-----------------------------------------------------------------------

The political struggle that eventually produced the Independent Treasury
began with Andrew Jackson's decision to destroy the Second Bank of the
United States (BUS), whose recharter he vetoed in 1832.  The Bank,
chartered in 1816 with a 20-year term, had served as the federal
government's fiscal agent: the Treasury deposited its revenues in the
BUS and its branches, which then acted as depositories and transfer
agents for the government's funds nationwide.

Jackson's veto message---drafted largely by Attorney General Roger Taney
and advisor Amos Kendall---condemned the BUS as an unconstitutional
monopoly, a corrupt engine of privilege, and a dangerous concentration
of financial power in private hands.  After his re-election in 1832
Jackson moved aggressively: in 1833--34 Secretary of the Treasury
William Duane refused to remove deposits from the BUS, so Jackson
replaced him with Taney, who then distributed the federal deposits to a
shifting roster of state-chartered banks---dubbed \emph{``pet banks''}
by Whig opponents.

The \textbf{pet-bank interlude (1833--1840)} demonstrated exactly the
problem that would motivate the Independent Treasury idea.  The
government's money sat in dozens of state banks across the country.
Those banks counted federal deposits as part of their reserves and
expanded credit on top of them.  When the Treasury needed to collect or
transfer funds, it depended entirely on the solvency and cooperation of
private institutions.  This arrangement blurred the line between public
revenue and private banking capital.

The Panic of 1837---triggered partly by Jacksonian hard-money policies
(the Specie Circular of 1836) and partly by a contraction of British
credit---caused most state banks to \textbf{suspend specie payments} in
May 1837.  Overnight, the federal deposits held in those banks became
inaccessible or nominally valueless, leaving the Treasury functionally
unable to meet its obligations.  This financial crisis gave urgency and
concrete shape to the proposal for a completely separate, self-contained
treasury.

%-----------------------------------------------------------------------
\section{The Sub-Treasury Proposal Takes Shape (1837--1840)}
%-----------------------------------------------------------------------

The idea of divorcing the Treasury entirely from the banking system had
been circulating in hard-money Democratic circles for several years, but
President Martin Van Buren made it the centrepiece of his
administration's response to the Panic.  In a special session of
Congress in September 1837, Van Buren proposed what he called the
\textbf{``Divorce Bill''} or \textbf{Sub-Treasury plan}: the federal
government would cease depositing its revenues in any bank---state or
national---and would instead keep all public money in its own strong
rooms, to be received and paid out exclusively in specie (gold and
silver coin, or notes fully convertible into specie).

The intellectual case rested on several arguments advanced by Van Buren,
Treasury Secretary Levi Woodbury, and congressional Democrats such as
Senator John C.\ Calhoun (who briefly allied with Van Buren on this
issue):

\begin{enumerate}
  \item \textbf{Separation of fiscal from monetary power.}  Public
    revenues should not serve as the reserve base for private credit
    expansion.  Government money was a public trust, not a subsidy to
    private banking.

  \item \textbf{Hard-money discipline.}  Requiring specie in and out
    would put downward pressure on bank-note inflation and fluctuating
    paper currencies.

  \item \textbf{Security of public funds.}  Government money held in
    government vaults could not be lost to bank failures or suspensions.

  \item \textbf{Constitutional propriety.}  Article~I, Section~9 of the
    Constitution required that ``no money shall be drawn from the
    Treasury but in consequence of appropriations made by law.'' Keeping
    funds in private banks made genuine accountability difficult.
\end{enumerate}

The Whig opposition---led by Henry Clay, Daniel Webster, and later
Millard Fillmore---fought the plan as deflationary, impractical, and a
blow to commercial credit.  The bill passed the Senate in 1838 but was
defeated in the House.  Van Buren reintroduced it, and the
\textbf{Independent Treasury Act of 1840} finally passed and was signed
into law.  The Whigs repealed it in 1841 after William Henry Harrison's
election.  The debate then simmered through the Tyler administration,
with Tyler vetoing two Whig bank bills in 1841, leaving the country
without any coherent fiscal system.

%-----------------------------------------------------------------------
\section{The Independent Treasury Act of 1846}
%-----------------------------------------------------------------------

When James K.\ Polk won the presidency in 1844 on an expansionist
Democratic platform, re-establishing the Independent Treasury was a
central fiscal priority.  Secretary of the Treasury Robert J.\ Walker
strongly supported it.  The \textbf{Independent Treasury Act of
August~6, 1846} restored and made permanent the system Van Buren had
briefly established.

The 1846 Act created a network of \textbf{sub-treasuries} (technically
called ``offices of the Assistant Treasurer of the United States'') at
key financial centres: New York, Boston, Philadelphia, Charleston, New
Orleans, St.\ Louis, and later San Francisco (after 1848).  The
Treasurer of the United States in Washington stood at the apex of the
system.

%-----------------------------------------------------------------------
\section{How the Independent Treasury Operated}
%-----------------------------------------------------------------------

\paragraph{Revenue collection.}
All federal revenues---customs duties, land sales, and other
receipts---had to be paid directly into the sub-treasuries and the
Treasury itself in \textbf{lawful money}: initially specie only, then
gradually also in Treasury notes and later in national bank notes after
Civil War-era amendments.

\paragraph{Custody.}
Federal funds were held physically in Treasury vaults.  There was no
commingling with private funds; no bank held government deposits.  Vault
holdings at each sub-treasury were reported regularly to Congress.

\paragraph{Disbursement.}
Federal expenditures---paying contractors, soldiers, pensioners,
bondholders---were made by warrant drawn on the Treasury or a
sub-treasury, payable in lawful money from those same vaults.

\paragraph{Transfer.}
The sub-treasury network allowed the government to move funds between
regions.  A warrant drawn in Washington could be presented at New York,
enabling payments across the country without the Treasury having to hold
a New York bank account.  This was important for customs receipts
(enormous in New York) and for disbursements in distant locations.

%-----------------------------------------------------------------------
\section{How the System Constrained the Secretary of the Treasury}
%-----------------------------------------------------------------------

The Independent Treasury Act severely curtailed the discretionary power
of the Secretary of the Treasury in several important respects.

\paragraph{No deposits in state banks.}
The entire premise of the Act was that the Secretary \emph{could not}
place federal revenue in state-chartered private banks, regardless of
how solvent or convenient they might be.  This reversed both the BUS era
(national fiscal agent) and the pet-bank era (discretionary executive
deposits).  The fiscal operations of the government were legally sealed
off from the state banking system.

\paragraph{No deposits in nationally chartered private banks.}
After the National Banking Acts of 1863--64 created a system of
federally chartered national banks, one might have expected the Treasury
to use them as depositories.  For ordinary Treasury operations under the
Independent Treasury framework, it did not: the sub-treasury vaults
remained the legal home of public funds.  Limited exceptions for
US-bond-secured deposits of national banks were introduced over time,
but these were narrow carve-outs, not a general depositary relationship.

\paragraph{Specie requirements as a monetary constraint.}
The early ``specie clause'' of the system meant that the Secretary's
receipts and payments had to be in coin (or approved equivalents).  This
removed the discretion to accept or issue bank paper on behalf of the
government, reducing the Treasury's ability to accommodate the banking
system during financial stringency.

\paragraph{Pro-cyclical fiscal mechanics.}
Because the Treasury held its own cash balances independently, revenue
surpluses physically drained specie from the banking system into
Treasury vaults, contracting bank reserves.  Conversely, government
expenditure injected specie back into circulation.  This made Treasury
cash management a significant macroeconomic variable---one the Secretary
had limited discretion to smooth, since depositing surpluses in banks to
cushion the contraction was precisely what the Act forbade.  This
pro-cyclical character became a chronic complaint about the system in
the late nineteenth century, most acutely during the panics of 1873,
1893, and 1907.

\paragraph{Long-term legacy and eventual repeal.}
The Independent Treasury survived the Civil War (with modifications,
since greenbacks and national bank notes were added to acceptable payment
media) and persisted in modified form until the \textbf{Federal Reserve
Act of 1913}, which finally transferred the government's fiscal agency
functions to the newly created Federal Reserve Banks.  The passage of
the Fed Act was partly motivated by the demonstrated failure of the
Independent Treasury to prevent the Panic of 1907---precisely because
the Secretary lacked the authority to freely deposit Treasury surpluses
in banks to relieve financial stringency without special emergency
legislation.

%-----------------------------------------------------------------------
\section{Eroding the Independent Treasury: Gage, Shaw, and Active
Treasury Management (1890s--1907)}
%-----------------------------------------------------------------------

Despite its apparently rigid statutory character, the Independent
Treasury framework was progressively hollowed out between the 1890s
and 1907 by a series of administrative innovations initiated by two
successive Secretaries of the Treasury: Lyman J.\ Gage (1897--1902)
and Leslie M.\ Shaw (1902--1907).  Both men believed that the Treasury
had an obligation---and the latent authority---to deploy government
funds actively in order to stabilise the money market.  Their policies,
which were controversial in their own day, are now chiefly remembered
as a failed proto-central-banking experiment whose limits helped prepare
the intellectual and political ground for the Federal Reserve Act of
1913.

%-----------------------------------------------------------------------
\subsection{Legal Foundations: National Bank Depositories}
%-----------------------------------------------------------------------

The Independent Treasury Act of 1846 prohibited deposits in state
banks, and the subsequent system was designed around specie held in
government vaults.  The National Banking Acts of 1863--64, however,
created a parallel channel through which government money could
re-enter the banking system: the Treasury was authorised to designate
certain nationally chartered banks as \textbf{government depositories},
provided those banks pledged United States bonds as collateral for the
deposits.  Customs receipts, in particular, could be placed in these
\emph{national bank depositories} rather than in sub-treasury vaults.
This was not a new general licence to deposit freely; the bond-security
requirement was meant to make it fiscally safe and limited in scope.
The Gold Standard Act of March 14, 1900---passed under Gage's
initiative---strengthened this channel further by authorising the
Secretary to use deposits in national banks as a mechanism for
maintaining the gold reserve.  It was these provisions that Gage and
Shaw would subsequently stretch far beyond their original intent.

%-----------------------------------------------------------------------
\subsection{Secretary Lyman Gage (1897--1902)}
%-----------------------------------------------------------------------

Gage arrived at the Treasury as president of the First National Bank of
Chicago, one of the most sophisticated financiers of his era.  He
presided over the country's recovery from the prolonged depression
following the Panic of 1893 and brought to his office a banker's
instinct that Treasury operations should be managed with an eye to
their monetary effects.

During episodes of autumnal tightness in the money market---when
interior banks withdrew deposits from New York to finance the harvest
and crop-moving season---Gage began the practice of \textbf{pre-emptive
deposit expansion}: placing government funds in New York national bank
depositories ahead of the seasonal stringency, thereby cushioning the
shock to bank reserves.  Conversely, when Treasury surpluses threatened
to drain reserves, he would accelerate bond purchases (effectively
returning cash to the market) or increase government deposits.  These
operations were conducted under the ``government depository'' authority
of the national banking legislation, but Gage interpreted that authority
broadly, increasing the number of designated depository banks well
beyond traditional norms and adjusting deposit levels with a deliberate
counter-cyclical logic.

Critics noted that these operations sat in considerable tension with the
spirit of the Independent Treasury Act, which had been premised on a
strict separation of government from banking.  Gage defended the
practice in his annual reports as a legitimate exercise of fiscal
management, not a violation of the law.  His predecessor's careful
compliance with the Independent Treasury's restrictive logic had, in
Gage's view, made the Treasury an instrument of monetary disruption
rather than stability.

%-----------------------------------------------------------------------
\subsection{Secretary Leslie M.\ Shaw (1902--1907)}
%-----------------------------------------------------------------------

Gage's successor, Leslie M.\ Shaw, took the same basic approach and
extended it dramatically, turning what Gage had done episodically into
a systematic programme of Treasury money-market management.  Shaw
served under Theodore Roosevelt and believed, with considerable
conviction, that the Secretary of the Treasury possessed---or ought to
possess---powers analogous to those of a central bank.

\paragraph{Expansion of depository banks and removal of reserve
requirements.}
In 1902, early in his tenure, Shaw announced that national banks
holding government deposits would no longer be required to maintain
cash reserves against those deposits.  Since national banks were
required by the National Banking Acts to hold a specified fraction of
their deposit liabilities as lawful-money reserves, this ruling
effectively allowed banks to treat public funds as costless additions
to their reserve base---a substantial subsidy to the banking system
and one that dramatically amplified the monetary effect of Treasury
deposits.  Shaw also sharply expanded the list of national banks
designated as government depositories, spreading the government's
money-market presence across the country.

\paragraph{Pre-emptive and remedial depositing.}
Like Gage, Shaw used government deposits to offset seasonal stringency
in the autumn.  But he went further: he deposited funds in advance,
in amounts calibrated to the expected demand for currency, and
withdrew them in spring when conditions eased.  He also used bond
purchases and deposit operations to respond to acute episodes of
financial stress, effectively acting as a lender of last resort
through the deposit channel rather than through any formal reserve
mechanism.

\paragraph{The 1906 Annual Report and the claim of central-bank
powers.}
Shaw's most celebrated---and controversial---statement of his
philosophy appeared in his \textbf{Annual Report for fiscal year
1906}, widely cited by contemporaries and by later historians.
Shaw there argued that the Secretary of the Treasury, if given
adequate funds and full discretion over their placement, could
effectively regulate money-market conditions throughout the country.
He maintained that the Treasury's deposit operations could serve all
the essential functions of a central bank's open-market and
discount-rate operations.  His specific claim was that with
``\$100{,}000{,}000 placed to the credit of the Treasurer in
national bank depositories, he could, at the option of the
Secretary, make money easy or tight at any time and in any
section''---a striking assertion of executive monetary power.
Shaw admitted that the legal authority for some of his practices was
uncertain, but argued that the public interest justified a broad
construction of the Secretary's inherent discretionary powers.

Milton Friedman and Anna Jacobson Schwartz quote from and discuss
this report in \emph{A Monetary History of the United States,
1867--1960} (1963), Chapter~5.  They use Shaw's statements as
evidence that the monetary instability of the pre-Federal Reserve
period was not simply the product of ignorance: the intellectual
tools for counter-cyclical monetary management were understood, and
at least one senior official had tried to apply them.  For Friedman
and Schwartz, however, Shaw's experience also illustrated the
decisive limitation of working within the Independent Treasury
framework---the Secretary's powers were ad hoc, legally fragile,
dependent on surplus revenue, and lacked the systematic authority
(to expand the money supply through discounting, for example) of a
genuine central bank.

%-----------------------------------------------------------------------
\subsection{Critical Reaction and the A.\ Piatt Andrew Assessment}
%-----------------------------------------------------------------------

Shaw's operations attracted immediate scholarly attention.  The most
thorough contemporary analysis was A.\ Piatt Andrew, ``The Treasury and
the Banks Under Secretary Shaw,'' \emph{Quarterly Journal of Economics}
21 (1907): 519--568.  Andrew documented Shaw's practices in careful
quantitative detail and was broadly critical: he argued that by
removing reserve requirements against public deposits, Shaw had
artificially inflated the effective reserve base of the banking system,
contributing to the credit expansion that preceded the Panic of 1907.
Andrew elaborated this argument in a follow-up piece, ``The United
States Treasury and the Money Market: The Partial Responsibility of
Secretaries Gage and Shaw for the Crisis of 1907,'' \emph{American
Economic Association Quarterly} 9 (1908): 218--231, which assigned
partial blame for the 1907 panic to the monetary expansion encouraged
by Treasury deposit policy.

A more sympathetic reading appeared in Richard H.\ Timberlake, ``Mr.\
Shaw and His Critics: Monetary Policy in the Golden Era Reviewed,''
\emph{Quarterly Journal of Economics} 77 (1963): 40--54, which argued
that Shaw's policies were a reasonable response to the structural
deficiency of an inelastic currency and that the criticism they
attracted was partly motivated by institutional rivalries.  Eugene B.\
Patton offered an intermediate legal-historical assessment in
``Secretary Shaw and Precedents as to Treasury Control over the Money
Market,'' \emph{Journal of Political Economy} 15 (1907): 65--87.

%-----------------------------------------------------------------------
\subsection{Significance and Legacy}
%-----------------------------------------------------------------------

The Gage--Shaw episode is significant for several reasons.

First, it revealed the irreducible tension within the Independent
Treasury system between fiscal neutrality and monetary stability.
A government that collects taxes and spends them cannot, in
a developed financial system, avoid having its cash management affect
the money market.  The 1846 Act's architects had hoped to insulate
public finance from banking; the late-nineteenth-century experience
showed that isolation cut both ways, making the Treasury itself a
source of destabilising shocks.

Second, the episode demonstrated that the formal constraints of the
Independent Treasury Act could be eroded through administrative
interpretation.  Gage and Shaw did not amend or repeal the Act; they
worked around it by exploiting the depository-bank provisions of the
national banking legislation, interpreting those provisions broadly and
removing conditions (bond security, reserve requirements) that had been
intended to limit their scope.

Third---and most consequentially for monetary history---the Gage--Shaw
experiment made visible both the possibility and the limits of
Treasury-based monetary management.  The Secretary's authority was
bounded by the size of the fiscal surplus (he could not create
reserves, only move existing ones), lacked the discount-rate mechanism
through which a central bank could influence the cost of credit
directly, and was politically vulnerable in ways that a quasi-independent
reserve bank was not.  These deficiencies were central to the arguments
advanced by the National Monetary Commission (1909--12) and ultimately
shaped the design of the Federal Reserve System created in 1913.

%-----------------------------------------------------------------------
\section{Summary Table}
%-----------------------------------------------------------------------

\begin{table}[ht]
\centering
\caption{Federal Fiscal Arrangements and Secretarial Discretion, 1791--1913}
\label{tab:fiscal_arrangements}
\renewcommand{\arraystretch}{1.4}
\begin{tabular}{>{\raggedright}p{3.2cm}
                >{\raggedright}p{4.8cm}
                >{\raggedright\arraybackslash}p{4.8cm}}
\toprule
\textbf{Phase} & \textbf{Arrangement} & \textbf{Secretary's Discretion} \\
\midrule
First BUS (1791--1811)
  & Federal deposits in BUS branches
  & Low---BUS was statutory fiscal agent \\
Interregnum (1811--1816)
  & State banks
  & High---informal and chaotic \\
Second BUS (1816--1836)
  & Federal deposits in BUS branches
  & Low---BUS was statutory fiscal agent \\
Pet banks (1833--1840)
  & Discretionary state bank deposits
  & Very high---executive choice \\
Independent Treasury (1840--41; 1846--1913)
  & Sub-treasury vaults only
  & Very low---banks statutorily excluded \\
Federal Reserve era (1914--present)
  & Fed as fiscal agent; later also commercial banks
  & Regulated but flexible \\
\bottomrule
\end{tabular}
\end{table}

The Independent Treasury thus represented the triumph of hard-money
Jacksonian ideology: the state was to be fiscally self-sufficient,
neither dependent on nor a subsidy to the private banking system.  The
price was a rigid, often deflationary fiscal mechanism that periodically
destabilized the very commercial economy it was meant to insulate the
government from.

%-----------------------------------------------------------------------
\section{Fiscal versus Monetary Dominance in Nineteenth-Century America:
Reflections on Sargent--Wallace}
%-----------------------------------------------------------------------

\subsection{The Sargent--Wallace Framework}

Thomas Sargent and Neil Wallace's 1981 paper ``Some Unpleasant
Monetarist Arithmetic'' (\emph{Federal Reserve Bank of Minneapolis
Quarterly Review}, Fall 1981) posed a question that the nineteenth-century
American fiscal experience illuminates with unusual sharpness: in a
world where a fiscal authority and a monetary authority make
independent decisions, which one \emph{dominates} in determining
the price level?

The core insight is a game-theoretic one.  The government's
consolidated budget constraint links the real primary fiscal surplus,
the stock of public debt, and the rate of money creation.  If the
fiscal authority commits first and credibly---fixing its deficit
path independently of monetary conditions---then the monetary
authority has no free choice: it must eventually generate enough
seigniorage to satisfy the government's intertemporal budget
constraint.  A monetary tightening today does not make inflation
disappear; it can only postpone it, and by raising the real interest
rate it may actually \emph{increase} the total future inflation needed
to retire the accumulated debt.  The resolution of this tension---
whether it takes the form of price-level determination by the fiscal
authority (the ``fiscal theory of the price level,'' as later
elaborated by Eric Leeper, Michael Woodford, and John Cochrane) or
an eventual monetary accommodation of fiscal needs---depends on
which authority operates under a binding constraint and which
operates with residual freedom.

The nineteenth-century United States provided no clean test of this
framework because it had no operationally independent monetary
authority for most of the period.  But that absence is itself
theoretically revealing: the institutional history can be read as a
series of experiments in creating, dissolving, and improvising
monetary authorities, each constrained in different ways by the
fiscal authority (Congress) that had delegated power to them.

%-----------------------------------------------------------------------
\subsection{The Monetary Constitution as Fiscal Legislation: Congress
as Ultimate Monetary Authority}
%-----------------------------------------------------------------------

In the United States, monetary institutions were created by
\emph{ordinary statute}, not by constitutional design.  Article~I,
Section~8 of the Constitution grants Congress the power ``to coin
money and regulate the value thereof,'' but it says nothing about
who should issue paper money, manage a central bank, or determine
the money supply.  As a result, every major monetary institution
of the nineteenth century was created, modified, and ultimately
destroyed by the legislative branch acting in its fiscal-political
interest.

This constitutional subordination of monetary to fiscal authority
is the deepest structural feature of nineteenth-century American
monetary history.  When Congress fixed the bimetallic ratio at
$15:1$ silver to gold in the Coinage Act of 1792 and then revised
it to $16:1$ in 1834, it was making a monetary policy decision in
the full modern sense: it set the relative price of the monetary
metals and thereby the long-run price level (given gold and silver
supply).  When Congress passed the Resumption Act of 1875 or the
Gold Standard Act of 1900, it was committing the monetary system to
a particular rule.  And when Congress issued an explicit quantity of
legal-tender greenbacks in 1862--63---\$450 million in non-interest-bearing
fiat currency---it was directly financing a fiscal deficit by money
creation, the textbook case of fiscal dominance in the Sargent--Wallace
sense.

The greenback episode (1862--79) is, in fact, the cleanest instance
of fiscal dominance in American history.  Congress determined the
quantity of money directly, without the intermediation of any monetary
authority.  The resulting inflation was the predictable arithmetic
consequence: the price level roughly doubled between 1861 and 1864.
Resumption of specie payments in January 1879 then required a
prolonged and painful deflation as the money stock was contracted to
restore the pre-war gold parity.  Sargent's analysis of ``the ends of
four big inflations'' (1982) applies to the greenback episode in
compressed form: the inflation ended because a credible fiscal-monetary
regime change---commitment to specie resumption backed by Congressional
action---altered expectations, with the price-level fall occurring
ahead of the formal resumption date.

Outside the greenback episode, the operative monetary commitment
was the gold standard.  One can interpret the gold standard as a device
that \emph{removed monetary policy from the hands of any human
authority}---Congress, BUS, or Secretary alike---and delegated it to
the international gold market.  Under a credible gold standard, the
domestic price level is anchored by the world price of gold, and no
domestic institution need (or indeed can) conduct independent monetary
policy.  In Leeper's terminology, the gold standard enforced
``passive'' monetary policy: the money supply adjusted endogenously
to maintain the gold parity, and no discretionary authority was needed
or permitted.  The fiscal authority was correspondingly constrained to
a ``Ricardian'' regime: it could not run persistent deficits without
eventually raising revenue, because there was no seigniorage valve
available to resolve unsustainable debt dynamics.

The institutional history surveyed in this document can therefore be
read as a chronicle not of \emph{who made monetary policy} in any
active discretionary sense, but of \emph{whose decisions could
undermine the gold-standard commitment}---and of institutional
arrangements designed to prevent that from happening.

%-----------------------------------------------------------------------
\subsection{The First and Second Banks as Delegated Monetary--Fiscal
Intermediaries}
%-----------------------------------------------------------------------

Alexander Hamilton's design for the First Bank of the United States
(1791) was a sophisticated solution to a fiscal--monetary coordination
problem that can be framed in modern terms.  The new republic's public
debt was unsustainable under the Articles of Confederation; no credible
commitment to repayment existed; tax revenues were inadequate and
unreliable.  Hamilton's plan---assumption of state debts, funding at par,
a national bank as fiscal agent---was in effect a scheme to move the
fiscal authority from a non-Ricardian to a Ricardian regime.  The
bank's role was not merely to provide a depository for Treasury funds;
it was to make the government's commitments credible by intertwining
the fiscal fortunes of the government with those of private
wealth-holders (bank shareholders and note-holders).

From a Sargent--Wallace perspective, the BUS was a device for
\emph{fiscal stabilisation} rather than monetary policy in the modern
sense.  Its monetary effects---the banknotes it issued, its influence
on the note-issuing behaviour of state banks---were largely
by-products of its fiscal functions.  Crucially, the BUS had no
statutory mandate to stabilise the price level or the money supply; it
had a mandate to serve as the government's fiscal agent and to
maintain convertibility of its own notes.

The Second Bank under Nicholas Biddle moved closer to something
resembling active monetary management.  Biddle used the BUS's position
as the nation's largest holder of state-bank notes to discipline state
banks: by accumulating and then presenting their notes for redemption
in specie, the BUS could force contractions on banks that had expanded
too aggressively.  This was a primitive open-market operation conducted
through the clearing mechanism rather than through bond purchases.
Biddle was explicitly aware of the monetary-policy dimension of these
actions and defended them publicly as stabilising.  But his authority
to conduct them rested entirely on the BUS's privately profitable
interest in sound banking, not on any statutory mandate.  Jackson's
destruction of the BUS eliminated whatever quasi-monetary-authority
function it had performed, leaving---as Section~\ref{sec:monetary_who}
hints---a vacuum that neither the pet-bank system nor the Independent
Treasury could fill.

%-----------------------------------------------------------------------
\subsection{The Secretary of the Treasury as De Facto Monetary
Authority}
\label{sec:monetary_who}
%-----------------------------------------------------------------------

With the destruction of the Second BUS, the United States was left
with a situation in which the \emph{Secretary of the Treasury}
was, in practice if not in law, the closest thing to a monetary
authority the country possessed---at least for managing the interface
between the government's finances and the banking system.

This was an authority derived entirely from \emph{delegation} by
Congress, and delegation of a highly imperfect kind.  The Secretary's
monetary-relevant powers consisted of:

\begin{itemize}
  \item The decision about \emph{timing and form of debt issuance}.
    When the Treasury needed to borrow---as it did heavily during wars
    and recessions---the structure of the bonds issued (maturity,
    coupon, tax-exempt status, eligibility as national-bank reserves
    under the National Banking Acts) had first-order effects on bank
    reserves and the money supply.  Congress authorised specific
    borrowing programmes, but the Secretary had substantial discretion
    over the terms of individual offerings.

  \item The decision about \emph{bond retirement}.  During the
    long period of postbellum fiscal surpluses (roughly 1866--93),
    the Treasury accumulated large cash balances that it deployed to
    buy back debt.  Each bond purchase injected money into the
    banking system; each period of inaction drained it.  This
    created a de facto open-market operation, conducted for purely
    fiscal reasons, but with monetary consequences that the Secretary
    could---and eventually did---manage deliberately.

  \item The \emph{deposit-placement} powers described in Sections~5
    and~6 of this document: under the pet-bank era and subsequently
    under the Gage--Shaw interpretations of the national-bank
    depository legislation.
\end{itemize}

What is striking from the perspective of Sargent--Wallace is that
in each of these cases the Secretary's ``monetary'' actions were
driven primarily by \emph{fiscal} logic.  The surplus that created
destabilising Treasury cash accumulations arose from tariff and land
revenues set by Congress for protectionist and revenue-raising
purposes, not from any monetary objective.  The seasonal tightness
that Gage and Shaw sought to offset arose from the calendar of tax
collection and crop-marketing, not from any monetary-policy design.
The bond-retirement operations of the postbellum era were driven by
the political imperatives of debt reduction, not by money-supply
targeting.

In Sargent--Wallace language: the fiscal authority (Congress) was
dominant, running surpluses or deficits as its own political economy
dictated, and the monetary consequences were residual---handled by
whoever happened to be Secretary rather than by any authority designed
for that purpose.  The ``monetary policy'' of the nineteenth-century
United States was, in this reading, almost entirely \emph{endogenous
to fiscal policy}, with the gold standard (where it held) providing
the only exogenous nominal anchor.

%-----------------------------------------------------------------------
\subsection{The Gage--Shaw Episode as Informal Proto--Central Banking
and its Theoretical Significance}
%-----------------------------------------------------------------------

When Secretaries Gage and Shaw began deliberately managing Treasury
deposit operations to stabilise the money market, they were groping
toward a solution to the problem that Sargent and Wallace would later
formalise.  The problem, as they understood it, was that the
Independent Treasury's rigid separation of government cash from the
banking system made the Treasury \emph{itself} a source of monetary
shocks: surpluses drained reserves, deficits replenished them, and
neither movement was calibrated to the needs of the money market.

Shaw's claim in his 1906 Annual Report---that with \$100~million in
discretionary deposits he could ``make money easy or tight at any
time and in any section''---is essentially a claim to possess a
proto--open-market--operation tool.  He was describing a policy that,
in modern terms, would be called \emph{monetary policy conducted by
the fiscal authority through its cash management}.  This is
precisely the configuration that Sargent and Wallace identify as the
problematic case: a single institution (the Treasury) in which fiscal
and monetary decisions are fused, with no independent authority
capable of credibly pre-committing to a monetary rule.

The limitations Shaw encountered are equally illuminating.  His
deposit operations could move existing reserves between the Treasury
and the banking system, but they could not \emph{create} reserves:
the total was bounded by the fiscal surplus.  He had no discount
window; he could not lend to banks directly.  His actions were legally
fragile---built on broad readings of limited statutory language that
could be challenged or reversed by the next administration.  And
because his authority derived from the executive branch rather than
from an independent institution, it was subject to political pressures
that a modern central bank is specifically designed to resist.

The Panic of 1907, which occurred just months after Shaw left office,
demonstrated these limits in a manner that A.~Piatt Andrew and others
publicised immediately.  The Treasury under Shaw's successor (George
Cortelyou, and then Franklin MacVeagh) once again used emergency
deposit operations to try to contain the panic, but lacked the
institutional capacity to act as a lender of last resort.  The
resolution of the panic required instead the private coordination of
J.~P.~Morgan---a private actor substituting for a public monetary
authority that did not yet exist.

%-----------------------------------------------------------------------
\subsection{Institutional Design Lessons: From the Independent Treasury
to the Federal Reserve}
%-----------------------------------------------------------------------

The history traced in this document suggests a specific progression in
the institutional management of the fiscal--monetary interface in
American history.

\begin{enumerate}

  \item \textbf{Congressional monetary supremacy (1787--1791).}
    The Constitution vested monetary power entirely in Congress, with
    no institutional mechanism for independent monetary management.
    This is pure fiscal dominance: monetary outcomes are whatever
    Congress, acting for fiscal-political reasons, dictates.

  \item \textbf{Hamiltonian delegation to a national bank (1791--1811,
    1816--1836).}  The BUS created a quasi-independent institution with
    both fiscal-agent and proto-monetary functions.  The bank's
    private-shareholder structure gave it incentives (profit,
    solvency) broadly aligned with monetary stability, but it had
    no price-level mandate and no lender-of-last-resort function.
    The delegation was revocable by Congress at any time, as Jackson
    demonstrated.

  \item \textbf{Fragmented delegation under the Independent Treasury
    (1840/46--1913).}  No monetary authority existed.  The gold
    standard provided a nominal anchor exogenous to any domestic
    institution, but could not prevent the pro-cyclical fiscal
    shocks described above.  The Secretary of the Treasury
    exercised residual quasi-monetary powers, but without statutory
    mandate, institutional independence, or adequate instruments.

  \item \textbf{The Gage--Shaw interlude (1897--1907).}  An informal,
    personalised, legally uncertain attempt to conduct monetary
    management through Treasury cash operations.  This was ``monetary
    policy by the fiscal authority''---precisely the configuration
    that Sargent--Wallace show to be problematic.

  \item \textbf{The Federal Reserve Act of 1913.}  The Federal Reserve
    represented a deliberate attempt to solve the fiscal--monetary
    coordination problem by \emph{separating} monetary from fiscal
    authority.  The Fed was given instruments (the discount rate,
    open-market operations, required reserves) that the Secretary
    lacked; it was insulated from direct Congressional appropriation;
    and it was designed to respond to monetary and financial
    conditions rather than to the Treasury's fiscal calendar.
    Whether one regards this separation as achieving genuine monetary
    independence is a question that continued to animate economic
    debate---notably during the Treasury--Federal Reserve Accord of
    1951, when the Fed reasserted its independence from Treasury
    pressure to peg interest rates---and that Sargent--Wallace's
    arithmetic showed could never be fully resolved as long as the
    consolidated government budget constraint remains binding.

\end{enumerate}

The nineteenth-century American experience thus provides a remarkable
historical laboratory for the abstract propositions of the
Sargent--Wallace framework.  It shows what a monetary system looks
like when there is \emph{no} operationally independent monetary
authority---when monetary outcomes are determined partly by the gold
standard's exogenous discipline, partly by Congressional legislation,
and partly by the discretionary and legally contested actions of a
single executive officer (the Secretary of the Treasury) who was
simultaneously a fiscal and a quasi-monetary agent.  The eventual
creation of the Federal Reserve can be understood not merely as a
banking reform but as an institutional solution to the
fiscal--monetary-dominance problem: an attempt to create, for the
first time in American history, a monetary authority that was
genuinely separate from---though never fully independent of---the
fiscal authority that had dominated it for more than a century.

\end{document}
